What remains to be done in order to prove Theorem \ref{maintheorem1} and what constitutes our main task is therefore to show that the generalized Robertson conditions are equivalent to the conditions~(\ref{Ricciform}) on the Ricci tensor, thus generalizing the classical result of Eisenhart \cite{Eis1934} to Painlev\'e metrics. In order to do this, we will show that
\begin{gather}\label{adaptedRicci}
R_{\jalpha \kbeta}=-\frac{3}{4}\partial_{j_{\alpha}}\partial_{k_{\beta}}\log\left[\frac{(\det S)^{n-2}}{\big(s^{11}\big)^{l_{1}}\cdots \big(s^{r1}\big)^{l_{r}}}\right] + \frac{1}{4}T_{\jalpha \kbeta},
\end{gather}
where
\begin{gather}
T_{\jalpha \kbeta} =(l_{\alpha}+l_{\beta}-2)\frac{s^{\alpha 1}s^{\beta 1}}{\det S}\partial_{j_{\alpha}}\partial_{k_{\beta}}\left(\frac{\det S}{s^{\alpha 1}s^{\beta 1}}\right)\nonumber\\
\hphantom{T_{\jalpha \kbeta} =}{} +\sum_{\gamma\neq \alpha,\beta=1}^{r}l_{\gamma}\left[\left(\partial_{j_{\alpha}}\log\frac{s^{\gamma 1}}{\det S}\right)\left(\partial_{k_{\beta}}\log\frac{s^{\alpha 1}}{\det S}\right)\right.\nonumber\\
\left.\hphantom{T_{\jalpha \kbeta} =}{}+\left(\partial_{j_{\alpha}}\log\frac{s^{\beta 1}}{\det S}\right)\left(\partial_{k_{\beta}}\log\frac{s^{\gamma 1}}{\det S}\right)
-\frac{\det S}{s^{\gamma 1}}\partial_{j_{\alpha}}\partial_{k_{\beta}}\left(\frac{s^{\gamma 1}}{\det S}\right)\right].\label{Tterms}
\end{gather}
We note that the expression (\ref{adaptedRicci}) of the off block-diagonal components of the Ricci tensor $R_{\jalpha \kbeta}$ is independent of the $r$ Riemannian block metrics $G_{\beta}$, $1\leq \beta \leq r$ defined by~(\ref{blockmetric}). We also remark that the first term in the expression~(\ref{Tterms}) of~$T_{\jalpha \kbeta}$, involving second derivatives, vanishes identically in the special case of St\"ackel metrics since the pre-factor $\lalpha+\lbeta-2$ is zero in that case.

Once we will have established (\ref{adaptedRicci}), it will then follow from Lemma~\ref{LeviCivita} and more precisely from the generalized Levi-Civita conditions~(\ref{genLeviCivita}) and~(\ref{secondorder}) that the generalized Robertson conditions (\ref{GenRobertson}) are indeed equivalent to the vanishing conditions~(\ref{Ricciform}) on the non-block diagonal components of the Ricci tensor. We therefore proceed to establish the form (\ref{adaptedRicci}) of the Ricci tensor for a Painlev\'e metric.

The expression of the Ricci tensor in terms of the Christoffel symbols is given by
\begin{gather}\label{Ricci}
R_{\jalpha \kbeta}=R^{l}_{{\phantom l} l \jalpha \kbeta}=\partial_{l}{\Gamma^{l}}_{\jalpha \kbeta}-\partial_{\jalpha}{\Gamma^{l}}_{l\kbeta}+{\Gamma^{m}}_{\jalpha \kbeta}{\Gamma^{l}}_{l m}-{\Gamma^{m}}_{l\kbeta}{\Gamma^{l}}_{\jalpha m},
\end{gather}
where the summation convention is applied with $1\leq l,m \leq n = \dim M$. In order to compute the right-hand side of (\ref{Ricci}), we will need expressions for the Christoffel symbols of a Painlev\'e metric (\ref{Painleveform}). Using the standard formulas
\begin{gather*}%\label{Christoffelgen}
\Gamma_{hji}=\frac{1}{2}(\partial_{h}g_{ji}+\partial_{j}g_{ih}-\partial_{i}g_{hj}), \qquad {\Gamma^{i}}_{hj}=g^{ik}\Gamma_{hjk},
\end{gather*}
and writing the Painlev\'e metric (\ref{Painleveform}) in block-diagonal form as
\[
{\rm d}s^{2}=\sum_{\alpha = 1}^{r}\sum_{\ialpha=1_{\alpha}}^{\lalpha}\sum_{\jalpha=1_{\alpha}}^{\lalpha}(g_{\alpha})_{\ialpha \jalpha}{\rm d}x^{\ialpha}{\rm d}x^{\jalpha},
\]
we obtain for fixed indices $1\leq \alpha,\beta \leq r$,
\begin{gather}
{\Gamma^{\ialpha}}_{\kalpha \jbeta}=\frac{1}{2}\sum_{\palpha=1_{\alpha}}^{\lalpha} \big(g^{\alpha}\big)^{\ialpha \palpha}\partial_{\jbeta}(g_{\alpha})_{\kalpha \palpha},\nonumber\\
 {\Gamma^{\jbeta}}_{\ialpha \kalpha}=-\frac{1}{2}\sum_{\pbeta=1_{\beta}}^{\lbeta}\big(g^{\beta}\big)^{\jbeta \pbeta}\partial_{\pbeta}(g_{\alpha})_{\ialpha \kalpha}\qquad \text{for} \quad \alpha \neq \beta,\label{Christoffel}
\end{gather}
and
\begin{gather}\label{Christoffelmixed}
{\Gamma^{\ialpha}}_{\halpha \kalpha}=\frac{1}{2}\sum_{\kalpha=1_{\alpha}}^{\lalpha} g^{\ialpha \kalpha}(\partial_{\halpha}g_{\jalpha \kalpha}+\partial_{\jalpha}g_{\halpha \kalpha}-\partial_{\kalpha}g_{\halpha \jalpha}).
\end{gather}
In view of the expressions (\ref{Christoffel}) of the Christoffel symbols, it is convenient to split the sum over~$l$ appearing in (\ref{Ricci}) into three sums, the first sum corresponding to the values of the summation index $l$ lying in groups of indices different from the groups corresponding to $\alpha$ and $\beta$, and the remaining two to the values of $l$ belonging to the groups of indices labelled by $\alpha$ and $\beta$ respectively. Thus we write
\begin{gather}\label{Riccidecomp}
R_{\jalpha \kbeta}=\sum_{l=1}^{n}R^{l}_{{\phantom l} l \jalpha \kbeta}=\sum_{\gamma \neq \alpha,\, \beta=1}^{r}\sum_{\pgamma=1_{\gamma}}^{\lgamma}R^{\pgamma}_{{\phantom l} \pgamma \jalpha \kbeta} + \sum_{\palpha =1_{\alpha}}^{\lalpha}R^{\palpha}_{{\phantom l} \palpha \jalpha \kbeta}+\sum_{\pbeta=1_{\beta}}^{\lbeta}R^{\pbeta}_{{\phantom l} \pbeta \jalpha \kbeta}.
\end{gather}
Let us begin with the first term. We have, for $\gamma\neq \alpha, \beta$,
\begin{gather}
\sum_{\pgamma = 1_{\gamma}}^{\lgamma}R^{\pgamma}_{{\phantom l} \pgamma \jalpha \kbeta}=-\sum_{\pgamma = 1_{\gamma}}^{\lgamma}\partial_{\jalpha}{\Gamma^{\pgamma}}_{\pgamma \kbeta} +\sum_{\pgamma = 1_{\gamma}}^{\lgamma}\sum_{\malpha=1_{\alpha}}^{\lalpha}{\Gamma^{\malpha}}_{\jalpha \kbeta}{\Gamma^{\pgamma}}_{\pgamma \malpha}\nonumber\\
\hphantom{\sum_{\pgamma = 1_{\gamma}}^{\lgamma}R^{\pgamma}_{{\phantom l} \pgamma \jalpha \kbeta}=}{}
+\sum_{\pgamma = 1_{\gamma}}^{\lgamma}\sum_{\mbeta=1_{\beta}}^{\lbeta}{\Gamma^{\mbeta}}_{\jalpha \kbeta}{\Gamma^{\pgamma}}_{\pgamma \mbeta} -\sum_{\pgamma = 1_{\gamma}}^{\lgamma}\sum_{\mgamma=1_{\gamma}}^{\lgamma}{\Gamma^{\mgamma}}_{\pgamma \kbeta}{\Gamma^{\pgamma}}_{\jalpha \mgamma},\label{Riemann1}
\end{gather}
where the summations have been written out explicitly to avoid notational ambiguities. We begin with evaluating the first-derivative term on the right-hand side of~(\ref{Riemann1}) using the expressions~(\ref{Christoffel}) for the Christoffel symbols, thus obtaining, for $\gamma \neq \alpha, \beta$,
\begin{gather}\label{derivgamma}
\sum_{\pgamma = 1_{\gamma}}^{\lgamma}\partial_{\jalpha}{\Gamma^{\pgamma}}_{\pgamma \kbeta} = \frac{1}{2}\sum_{\pgamma = 1_{\gamma}}^{\lgamma}\sum_{\ngamma = 1_{\gamma}}^{\lgamma}\partial_{\jalpha}\big(\big(g^{\gamma}\big)^{\pgamma \ngamma}\partial_{\kbeta}((g_{\gamma})_{\pgamma \ngamma})\big)=\frac{1}{2}\partial_{\jalpha}\partial_{\kbeta}\log(|g_{\gamma}|),
\end{gather}
where $|g_{\gamma}|:=\det(g_{\igamma \jgamma})$. It follows that
\[
\sum_{\gamma \neq \alpha,\, \beta=1}^{r}\sum_{\pgamma=1_{\gamma}}^{\lgamma}\partial_{\jalpha}{\Gamma^{\pgamma}}_{\pgamma \kbeta}= \frac{1}{2}\partial_{\jalpha}\partial_{\kbeta}\left(\prod_{\gamma \neq \alpha, \,\beta=1}^{r}\log(|g_{\gamma}|)\right).
\]
From the Painlev\'e form (\ref{Painleveform}), we have
\begin{gather}\label{Painleveblock}
g_{\igamma \jgamma}=\frac{\det S}{s^{\gamma 1}}(G_{\gamma})_{\igamma \jgamma},
\end{gather}
so that
\[
|g_{\gamma}|=\frac{(\det S)^{\lgamma}}{\big({s^{\gamma 1}}\big)^{\lgamma}}|G_{\gamma}|.
\]
Using the fact that $|G_{\gamma}|$ is a function of the variables ${\bf x}^{\gamma}$ only, we obtain
\begin{gather*}%\label{secondderiv1}
\sum_{\gamma \neq \alpha, \,\beta=1}^{r}\sum_{\pgamma=1_{\gamma}}^{\lgamma} \partial_{\jalpha}{\Gamma^{\pgamma}}_{\pgamma \kbeta}=\frac{1}{2}\sum_{\gamma \neq \alpha, \,\beta=1}^{r}\lgamma \partial_{\jalpha}\partial_{\kbeta}\left(\log\left(\frac{\det S}{s^{\gamma 1}}\right)\right).
\end{gather*}
We notice that the above expression is independent of the quadratic differential forms $G_{\alpha}$ defined by~(\ref{blockmetric}).

Next we evaluate the terms quadratic in the Christoffel symbols in the right-hand side of~(\ref{Riemann1}). Again, we use the fact that $\gamma\neq \alpha, \beta$ and the fact that in the Painlev\'e form~(\ref{Painleveform}), each of the quadratic differential forms $G_{\alpha}$ defined by~(\ref{blockmetric}) depends on the group of variables ${\bf x}^{\alpha}$ only. We have
\begin{gather}
\sum_{\malpha=1_{\alpha}}^{\lalpha} \sum_{\pgamma = 1_{\gamma}}^{\lgamma}{\Gamma^{\malpha}}_{\jalpha \kbeta}{\Gamma^{\pgamma}}_{\pgamma \malpha}\nonumber\\
\qquad{} =\frac{1}{4}\sum_{\malpha = 1_{\alpha}}^{l_{\alpha}} \sum_{\pgamma = 1_{\gamma}}^{\lgamma}\sum_{\nalpha=1_{\alpha}}^{l_{\alpha}}(g^{\alpha})^{\malpha \nalpha}\partial_{\kbeta}((g_{\alpha})_{\jalpha \nalpha})g^{\pgamma \hgamma}\partial_{\malpha}((g_{\gamma})_{\pgamma \hgamma})\nonumber\\
\qquad{} =\frac{1}{4}\sum_{\malpha=1_{\alpha}}^{l_{\alpha}}\sum_{\nalpha=1_{\alpha}}^{l_{\alpha}}(g^{\alpha})^{\malpha \nalpha}\partial_{\kbeta}((g_{\alpha})_{\jalpha \nalpha})\partial_{\malpha}\log(|g_{\gamma}|)\nonumber\\
\qquad{} =\frac{1}{4}\sum_{\malpha=1_{\alpha}}^{l_{\alpha}}\sum_{\nalpha=1_{\alpha}}^{l_{\alpha}}\frac{s^{\alpha 1}}{\det S}\big(G^{\alpha}\big)^{\malpha \nalpha}\partial_{\kbeta}\left(\frac{\det S}{s^{\alpha 1}}(G_{\alpha})_{\jalpha \nalpha}\right)\partial_{\malpha}\log(|g_{\gamma}|)\nonumber\\
\qquad{} =\frac{1}{4}\sum_{\malpha=1_{\alpha}}^{l_{\alpha}}\sum_{\nalpha=1_{\alpha}}^{l_{\alpha}}{\delta^{\malpha}}_{\jalpha}\partial_{\kbeta}\log\left(\frac{\det S}{s^{\alpha 1}}\right)\partial_{\malpha}\log \left(\frac{(\det S)^{\lgamma}}{(s^{\gamma 1})^{\lgamma}}\right)\nonumber\\
\qquad{} =\frac{1}{4}\left(\partial_{\kbeta}\log\left(\frac{s^{\alpha 1}}{\det S}\right)\right)\left(\partial_{\jalpha}\log \left(\frac{\big(s^{\gamma 1}\big)^{\lgamma}}{(\det S)^{\lgamma}}\right)\right).\label{quadr1}
\end{gather}
We notice that the above expression is again independent of the quadratic differential forms $G_{\alpha}$ defined by (\ref{blockmetric}). Likewise, we obtain for the next quadratic term in the Christoffel symbols that appears in the right-hand side of~(\ref{Riemann1}),
\begin{gather}\label{quadr2}
\sum_{\mbeta=1_{\beta}}^{\lbeta} \sum_{\pgamma = 1_{\gamma}}^{\lgamma}{\Gamma^{\mbeta}}_{\jalpha \kbeta}{\Gamma^{\pgamma}}_{\pgamma \mbeta}=\frac{1}{4}\left(\partial_{\jalpha}\log\left(\frac{s^{\beta 1}}{\det S}\right)\right)\left(\partial_{\kbeta}\log \left(\frac{\big(s^{\gamma 1}\big)^{\pgamma}}{(\det S)^{\pgamma}}\right)\right).
\end{gather}
For the third and final quadratic term, we have
\begin{align}
\sum_{\mgamma=1_{\gamma}}^{\lgamma}{\Gamma^{\mgamma}}_{\pgamma \kbeta}{\Gamma^{\pgamma}}_{\jalpha \mgamma} &=\frac{1}{4}\sum_{\mgamma=1_{\gamma}}^{\lgamma}{\delta^{\mgamma}}_{\pgamma}\left(\partial_{\kbeta}\log\left(\frac{s^{\gamma 1}}{\det S}\right)\right){\delta^{\pgamma}}_{\mgamma}\left(\partial_{\jalpha}\log\left(\frac{s^{\gamma 1}}{\det S}\right)\right)\nonumber\\
&=\frac{1}{4}\lgamma \left(\partial_{\kbeta}\log\left(\frac{s^{\gamma 1}}{\det S}\right)\right)\left(\partial_{\jalpha}\log\left(\frac{s^{\gamma 1}}{\det S}\right)\right),\label{quadr3}
\end{align}
which is likewise independent of the quadratic differential forms $G_{\alpha}$ defined by (\ref{blockmetric}). Putting together the expressions (\ref{derivgamma}), (\ref{quadr1}), (\ref{quadr2}) and (\ref{quadr3}) we obtain
\begin{gather}
\sum_{1\leq \gamma \neq \alpha, \beta\leq r}R^{\pgamma}_{{\phantom l} \pgamma \jalpha \kbeta}=\sum_{\gamma \neq \alpha, \,\beta =1}^{r}l_{\gamma}\left[\frac{1}{2}\partial_{\jalpha}\partial_{\kbeta}\log\left(\frac{s^{\gamma 1}}{\det S}\right)\right.\nonumber\\
\hphantom{\sum_{1\leq \gamma \neq \alpha, \beta\leq r}R^{\pgamma}_{{\phantom l} \pgamma \jalpha \kbeta}=}{}
+\frac{1}{4}\left(\partial_{\kbeta}\log\left(\frac{s^{\alpha 1}}{\det S}\right)\right)\left(\partial_{\jalpha}\log\left(\frac{s^{\gamma 1}}{\det S}\right)\right)\nonumber\\
\hphantom{\sum_{1\leq \gamma \neq \alpha, \beta\leq r}R^{\pgamma}_{{\phantom l} \pgamma \jalpha \kbeta}=}{}
+ \frac{1}{4}\left(\partial_{\jalpha}\log\left(\frac{s^{\beta 1}}{\det S}\right)\right)\left(\partial_{\kbeta}\log\left(\frac{s^{\gamma 1}}{\det S}\right)\right)\nonumber\\
\left.\hphantom{\sum_{1\leq \gamma \neq \alpha, \beta\leq r}R^{\pgamma}_{{\phantom l} \pgamma \jalpha \kbeta}=}{}
-\frac{1}{4}\left(\partial_{\kbeta}\log\left(\frac{s^{\gamma 1}}{\det S}\right)\right)\left(\partial_{\jalpha}\log\left(\frac{s^{\gamma 1}}{\det S}\right)\right)\right].\label{Riccidecomp1}
\end{gather}
We still need to evaluate the curvature components $R^{\palpha}_{{\phantom l} \palpha \jalpha \kbeta}$ and $R^{\pbeta}_{{\phantom l} \pbeta \jalpha \kbeta}$, which will require a~separate calculation. We have
\begin{gather}
\sum_{\palpha=1_{\alpha}}^{\lalpha}R^{\palpha}_{{\phantom l} \palpha \jalpha \kbeta} =\sum_{\palpha=1_{\alpha}}^{\lalpha}\partial_{\palpha}{\Gamma^{\palpha}}_{\jalpha \kbeta}-\sum_{\palpha=1_{\alpha}}^{\lalpha}\partial_{\jalpha}{\Gamma^{\palpha}}_{\palpha \kbeta}\nonumber\\
\hphantom{\sum_{\palpha=1_{\alpha}}^{\lalpha}R^{\palpha}_{{\phantom l} \palpha \jalpha \kbeta} =}{}
+\sum_{\palpha=1_{\alpha}}^{\lalpha}\sum_{\malpha = 1_{\alpha}}^{\lalpha}{\Gamma^{\malpha}}_{\jalpha \kbeta}{\Gamma^{\palpha}}_{\palpha \malpha}+\sum_{\palpha=1_{\alpha}}^{\lalpha}\sum_{\mbeta=1_{\beta}}^{\lbeta}{\Gamma^{\mbeta}}_{\jalpha \kbeta}{\Gamma^{\palpha}}_{\palpha \mbeta}\nonumber\\
\hphantom{\sum_{\palpha=1_{\alpha}}^{\lalpha}R^{\palpha}_{{\phantom l} \palpha \jalpha \kbeta} =}{}-\sum_{\palpha=1_{\alpha}}^{\lalpha}\sum_{\malpha=1_{\alpha}}^{\lalpha}{\Gamma^{\malpha}}_{\palpha \kbeta}{\Gamma^{\palpha}}_{\jalpha \malpha}-\sum_{\palpha=1_{\alpha}}^{\lalpha}\sum_{\mbeta=1_{\beta}}^{\lbeta}{\Gamma^{\mbeta}}_{\lalpha \kbeta}{\Gamma^{\palpha}}_{\jalpha \mbeta},\label{Riemann2}
\end{gather}
where again we have written out the summation signs explicitly to avoid notational ambiguities. We have, using (\ref{Christoffel}),
\begin{gather}\label{moreChristoffel}
{\Gamma^{\palpha}}_{\jalpha \kbeta}=\frac{1}{2}{\delta^{\palpha}}_{\jalpha}\partial_{\kbeta}\left(\log\left(\frac{\det S}{s^{\alpha 1}}\right)\right),
\end{gather}
so using the fact that the cofactor $s^{\alpha 1}$ is independent of ${\bf x}^{\alpha}$, we obtain
\begin{gather}\label{newderiv1}
\sum_{\palpha=1_{\alpha}}^{\lalpha}\partial_{\palpha}{\Gamma^{\palpha}}_{\jalpha \kbeta}=\frac{1}{2}\partial_{\jalpha}\partial_{\kbeta}\big(\log(\det S)\big),
\end{gather}
and a similar calculation gives
\begin{gather}\label{newderiv2}
\sum_{\palpha=1_{\alpha}}^{\lalpha}\partial_{\jalpha}{\Gamma^{\palpha}}_{\palpha \kbeta}=\frac{1}{2}\partial_{\jalpha}\partial_{\kbeta}\big(\log\big((\det S)^{\lalpha}\big)\big).
\end{gather}
We now evaluate the quadratic terms in the Christoffel symbols that appear in the curvature component~(\ref{Riemann2}). In order to do so, we substitute into the expression (\ref{Christoffelmixed}) of the Christoffel symbols the expressions
\[
(g_{\alpha})_{\ialpha \jalpha}=\frac{\det S}{s^{\alpha 1}}G_{\alpha},
\]
which result from the Painlev\'e form (\ref{Painleveform}). We obtain the following expressions for the Christoffel symbols,
\begin{gather}
{\Gamma^{\ialpha}}_{\halpha \jalpha}={\gamma^{\ialpha}}_{\halpha \jalpha} +\frac{1}{2}\Bigg[ {\delta^{\ialpha}}_{\jalpha}\partial_{\halpha}\big(\log{(\det S)}\big)+{\delta^{\ialpha}}_{\halpha}\partial_{\jalpha}\big(\log{(\det S)}\big)\nonumber\\
\hphantom{{\Gamma^{\ialpha}}_{\halpha \jalpha}=}{}-\sum_{\kalpha = 1_{\alpha}}^{\lalpha}\big(G^{\alpha}\big)^{\ialpha \kalpha}(G_{\alpha})_{\halpha \jalpha}\partial_{\kalpha}\big(\log{(\det S)}\big)\Bigg],\label{longChristoffel}
\end{gather}
where the ${\gamma^{\ialpha}}_{\halpha \jalpha}$ denote the Christoffel symbols of the block metric~$G_{\alpha}$ given by~(\ref{blockmetric}). It then follows from~(\ref{longChristoffel}) and~(\ref{moreChristoffel}) that
\begin{gather}
\sum_{\palpha=1_{\alpha}}^{\lalpha}\sum_{\malpha = 1_{\alpha}}^{\lalpha}{\Gamma^{\malpha}}_{\jalpha \kbeta}{\Gamma^{\palpha}}_{\palpha \malpha}\nonumber\\
\qquad{} =\frac{1}{2}\partial_{\kbeta}\left(\log\left(\frac{\det S}{s^{\alpha 1}}\right)\right)\left[\sum_{\palpha=1_{\alpha}}^{\lalpha}{\gamma^{\palpha}}_{\palpha \jalpha}+\frac{1}{2}\lalpha \partial_{\malpha}\big(\log{(\det S)}\big)\right].\label{newquadr1}
\end{gather}
Similarly, using (\ref{moreChristoffel}) and (\ref{longChristoffel}), we obtain
\begin{gather}
\sum_{\palpha=1_{\alpha}}^{\lalpha}\sum_{\malpha=1_{\alpha}}^{\lalpha}{\Gamma^{\malpha}}_{\palpha \kbeta}{\Gamma^{\palpha}}_{\jalpha \malpha}\nonumber\\
\qquad{} =\frac{1}{2}\partial_{\kbeta}\left(\log\left(\frac{\det S}{s^{\alpha 1}}\right)\right)\left[\sum_{\palpha=1_{\alpha}}^{\lalpha}{\gamma^{\palpha}}_{\jalpha \palpha}+\frac{1}{2}\lalpha \partial_{\malpha}\big(\log{(\det S)}\big)\right].\label{newquadr2}
\end{gather}
It follows therefore from (\ref{newquadr1}) and (\ref{newquadr2}) that
\[
\sum_{\palpha=1_{\alpha}}^{\lalpha}\sum_{\malpha = 1_{\alpha}}^{\lalpha}{\Gamma^{\malpha}}_{\jalpha \kbeta}{\Gamma^{\palpha}}_{\palpha \malpha}-\sum_{\palpha=1_{\alpha}}^{\lalpha}\sum_{\malpha=1_{\alpha}}^{\lalpha}{\Gamma^{\malpha}}_{\palpha \kbeta}{\Gamma^{\palpha}}_{\jalpha \malpha}=0.
\]
We now evaluate the remaining difference of two double sums in the expression (\ref{Riemann2}) of the curvature, using the expressions
\[
{\Gamma^{\mbeta}}_{\jalpha \kbeta}=\frac{1}{2}{\delta^{\mbeta}}_{\kbeta}\partial_{\jalpha}\left(\log\left(\frac{\det S}{s^{\beta 1}}  \right)\right), \qquad {\Gamma^{\mbeta}}_{\jalpha \kbeta} = \frac{1}{2}\partial_{\mbeta}\left(\log\left(\frac{(\det S)^{\lalpha}}{\big(s^{\beta 1}\big)^{\lalpha}} \right)\right),
\]
for the Christoffel symbols, to obtain
\begin{gather}
\sum_{\palpha=1_{\alpha}}^{\lalpha}\sum_{\malpha = 1_{\alpha}}^{\lalpha}{\Gamma^{\malpha}}_{\jalpha \kbeta}{\Gamma^{\palpha}}_{\palpha \malpha}-\sum_{\palpha=1_{\alpha}}^{\lalpha}\sum_{\mbeta=1_{\beta}}^{\lbeta}{\Gamma^{\mbeta}}_{\lalpha \kbeta}{\Gamma^{\palpha}}_{\jalpha \mbeta}\nonumber\\
\qquad{} =\frac{1}{4}(\lalpha-1)\partial_{\jalpha}\left(\log \left(\frac{\det S}{s^{\beta 1}}\right)\right)\partial_{\kbeta}\left(\log\left(\frac{\det S}{s^{\alpha 1}} \right)\right).\label{newquadr3}
\end{gather}
Substituting the expressions (\ref{newderiv1}), (\ref{newderiv2}), (\ref{newquadr1}), (\ref{newquadr2}) and (\ref{newquadr3}) into (\ref{Riemann2}), we obtain
\begin{gather}
\sum_{\palpha=1_{\alpha}}^{\lalpha}R^{\palpha}_{{\phantom l} \palpha \jalpha \kbeta}= \frac{1-\lalpha}{2} \partial_{\jalpha}\partial_{\kbeta} \big(\log(\det S)\big)\nonumber\\
\hphantom{\sum_{\palpha=1_{\alpha}}^{\lalpha}R^{\palpha}_{{\phantom l} \palpha \jalpha \kbeta}=}{}
+\frac{1}{4}(\lalpha-1)\partial_{\jalpha}\left(\log\left(\frac{\det S}{s^{\beta 1}}\right)\right)\partial_{\kbeta}\left(\log\left(\frac{\det S}{s^{\alpha 1}}\right)\right),\label{Riccidecomp2}
\end{gather}
and similarly
\begin{gather}\label{Riccidecomp3}
\sum_{\pbeta=1_{\beta}}^{\lalpha}R^{\pbeta}_{{\phantom l} \pbeta \jalpha \kbeta}=\frac{1-\lbeta}{2}\partial_{\jalpha}\partial_{\kbeta} \big(\log(\det S)\big)\nonumber\\
\hphantom{\sum_{\pbeta=1_{\beta}}^{\lalpha}R^{\pbeta}_{{\phantom l} \pbeta \jalpha \kbeta}=}{}
+\frac{1}{4}(\lbeta-1)\partial_{\jalpha}\left(\log\left(\frac{\det S}{s^{\beta 1}}\right)\right)\partial_{\kbeta}\left(\log\left(\frac{\det S}{s^{\alpha 1}}\right)\right).
\end{gather}
We now substitute the expressions (\ref{Riccidecomp1}), (\ref{Riccidecomp2}) and (\ref{Riccidecomp3}) into the decomposition~(\ref{Riccidecomp}) of the off-block diagonal Ricci curvature components $R_{\jalpha \kbeta}$ to obtain, for $\alpha \neq \beta$,
\begin{gather}
R_{\jalpha \kbeta}=-\frac{1}{2}\sum_{\gamma =1}^{r}\lgamma \partial_{\jalpha}\partial_{\kbeta}\left(\log\left(\frac{\det S}{s^{\gamma 1}}\right)\right)+\partial_{\jalpha}\partial_{\kbeta}\big(\log(\det S)\big)\nonumber\\
\hphantom{R_{\jalpha \kbeta}=}{} + \frac{1}{4}(\lalpha+\lbeta-2) \partial_{\jalpha}\left(\log\left(\frac{\det S}{s^{\beta 1}}\right)\right)\partial_{\kbeta}\left(\log\left(\frac{\det S}{s^{\alpha 1}}\right)\right)\nonumber\\
\hphantom{R_{\jalpha \kbeta}=}{}+\frac{1}{4}\sum_{\gamma \neq \alpha, \,\beta=1}^{r}\lgamma \left[\partial_{\kbeta}\left(\log\left(\frac{\det S}{s^{\alpha 1}}\right)\right) \partial_{\jalpha}\left(\log\left(\frac{\det S}{s^{\gamma 1}}\right)\right)\right.\nonumber\\
\hphantom{R_{\jalpha \kbeta}=}{}
+\partial_{\jalpha}\left(\log\left(\frac{\det S}{s^{\beta 1}}\right)\right)\partial_{\kbeta}\left(\log\left(\frac{\det S}{s^{\gamma 1}}\right)\right)\nonumber\\
\left.\hphantom{R_{\jalpha \kbeta}=}{} -\partial_{\kbeta}\left(\log\left(\frac{\det S}{s^{\gamma 1}}\right)\right)\partial_{\jalpha}\left(\log\left(\frac{\det S}{s^{\gamma 1}}\right)\right)\right].\label{Riccialmostfinal}
\end{gather}
We observe that since the cofactors $s^{\alpha 1}$ and $s^{\beta 1}$ are independent of the groups of variables ${\bf x}^{\alpha}$ and ${\bf x}^{\beta}$ respectively and since $\sum\limits_{\alpha=1}^{r}\lalpha = n$, we have
\begin{gather*}
-\frac{1}{2}\sum_{\gamma =1}^{r}\lgamma \partial_{\jalpha}\partial_{\kbeta}\left(\log\left(\frac{\det S}{s^{\gamma 1}}\right)\right)+\partial_{\jalpha}\partial_{\kbeta}\big(\log(\det S)\big)\\
\qquad{} =-\frac{1}{2}\partial_{j_{\alpha}}\partial_{k_{\beta}}\log\left[\frac{(\det S)^{n-2}}{\big(s^{11}\big)^{l_{1}}\cdots \big(s^{r1}\big)^{l_{r}}}\right].
\end{gather*}
We write
\begin{gather}
-\frac{1}{2}\partial_{j_{\alpha}}\partial_{k_{\beta}}\log\left[\frac{(\det S)^{n-2}}{\big(s^{11}\big)^{l_{1}}\cdots \big(s^{r1}\big)^{l_{r}}}\right]\nonumber\\
\qquad{} =-\frac{3}{4}\partial_{j_{\alpha}}\partial_{k_{\beta}}\log\left[\frac{(\det S)^{n-2}}{\big(s^{11}\big)^{l_{1}}\cdots \big(s^{r1}\big)^{l_{r}}}\right] +\frac{1}{4}\partial_{j_{\alpha}}\partial_{k_{\beta}}\log\left[\frac{(\det S)^{n-2}}{\big(s^{11}\big)^{l_{1}}\cdots \big(s^{r1}\big)^{l_{r}}}\right],\label{simplif}
\end{gather}
and compute the second term in (\ref{simplif}) as follows
\begin{gather*}
\frac{1}{4}\partial_{j_{\alpha}}\partial_{k_{\beta}}\log\left[\frac{(\det S)^{n-2}}{\big(s^{11}\big)^{l_{1}}\cdots \big(s^{r1}\big)^{l_{r}}}\right]\\
\qquad{} =\frac{1}{4}\partial_{j_{\alpha}}\partial_{k_{\beta}}\left[ \sum_{1\leq \gamma \neq \alpha, \beta\leq r} \log\left( \frac{(\det S)^{\lgamma}}{\big(s^{\gamma 1}\big)^{\lgamma}}\right)+\log\left(\frac{(\det S)^{\lalpha+\lbeta-2}}{s^{\alpha 1}s^{\beta 1}}\right)\right].
\end{gather*}
Evaluating the derivatives an using the fact that the cofactors $s^{\alpha 1}$ and $s^{\beta 1}$ are independent of the groups of variables ${\bf x}^{\alpha}$ and ${\bf x}^{\alpha}$, we conclude that the expression (\ref{Riccialmostfinal}) of the off-block diagonal Ricci curvature can be written as
\begin{gather*}
R_{\jalpha \kbeta} =-\frac{3}{4}\partial_{j_{\alpha}}\partial_{k_{\beta}}\log\left[\frac{(\det S)^{n-2}}{\big(s^{11}\big)^{l_{1}}\cdots \big(s^{r1}\big)^{l_{r}}}\right] + \frac{1}{4} (l_{\alpha}+l_{\beta}-2)\frac{s^{\alpha 1}s^{\beta 1}}{\det S}\partial_{j_{\alpha}}\partial_{k_{\beta}}\left(\frac{\det S}{s^{\alpha 1}s^{\beta 1}}\right)\\
\hphantom{R_{\jalpha \kbeta} =}{}+
\sum_{\gamma\neq \alpha,\beta=1}^{r}l_{\gamma}\left(\partial_{j_{\alpha}}\log\frac{s^{\gamma 1}}{\det S}\right)\left(\partial_{k_{\beta}}\log\frac{s^{\alpha 1}}{\det S}\right)
+\left(\partial_{j_{\alpha}}\log\frac{s^{\beta 1}}{\det S}\right)\left(\partial_{k_{\beta}}\log\frac{s^{\gamma 1}}{\det S}\right)\\
\hphantom{R_{\jalpha \kbeta} =}{}
-\frac{\det S}{s^{\gamma 1}}\partial_{\jalpha}\partial_{\kbeta}\left(\frac{s^{\gamma 1}}{\det S}\right).
\end{gather*}
This completes the proof of Theorem \ref{maintheorem1}.
